\section{Related Work}
\label{sec:related_work}

\subsection{Cross-modal music representation and correspondence learning}

A first line of multimodal music research arose from the need to connect heterogeneous music representations rather than to analyze performance behavior directly. Broad surveys by Simonetta et al.\ and M{\"u}ller et al.\ describe this area across audio, text, images, scores, metadata, and video, and emphasize that cross-modal retrieval and shared representation learning became early anchors for multimodal music research \citep{Simonetta2019Survey,Muller2019CrossModalSurvey}. Representative works include multimodal genre classification from audio, text, and visual cues \citep{Oramas2018Genre}, supervised audio-visual embedding for music video retrieval \citep{Zeng2018AVMusicRetrieval}, artistic correspondence learning between music and video \citep{Suris2022ArtisticCorrespondence}, and music-specific audio-language pretraining objectives such as CALL and MuLaP \citep{Manco2022CALLMusic,Manco2022MuLaP}. These methods established reusable objectives for aligning music modalities, but their dominant targets remain retrieval, tagging, or semantic matching rather than fast clip-level instrument decisions under local mismatch and low-latency constraints.

\subsection{Audio-visual music understanding and reasoning-oriented models}

Once shared representations became more reliable, a natural next step was to formulate music-centered multimodal understanding as a reasoning problem over dynamic audio-visual scenes. MUSIC-AVQA is representative of this shift: it casts audio-visual understanding as question answering and therefore evaluates whether a model can combine what is seen and what is heard at a finer semantic level than simple correspondence \citep{Li2022MusicAVQA}. Subsequent work specialized this idea for music by learning performance-aware representations for question answering \citep{Diao2024MPQA}, examining why general multimodal LLMs remain weak on music-focused AVQA \citep{You2025Complexity}, and introducing broader multimodal music understanding frameworks such as DeepResonance and MusiXQA together with explanation-oriented tools such as MusicLIME \citep{Mao2025DeepResonance,Chen2025MusiXQA,Sotirou2025MusicLIME}. Beyond music, video recognition studies based on spatial-temporal visual features also show that dynamic scene understanding benefits from explicit temporal modeling even when the target domain is not performance analysis \citep{Divya2025ActivityDetection}. These models substantially extend what can be asked of a multimodal music system. For the present problem, however, their objectives are broader than clip-level instrument classification, and their evaluations emphasize reasoning quality rather than whether a model can operate on short public performance clips under limited computational budget. They also do not isolate how local asynchrony, missing modalities, or camera degradation affect a lightweight decision pipeline.

\subsection{Performance-centric datasets and motion-aware audio-visual analysis}

A third line of work centers the performance event by collecting datasets with visible playing actions or by explicitly exploiting motion cues in audio-visual analysis. URMP provides multitrack ensemble recordings with synchronized audio, video, MIDI, and note annotations, and remains one of the most useful public references for multimodal performance study \citep{Li2019URMP}. Solos moves closer to per-performer visual behavior by offering solo videos together with frame-level skeletons for audio-visual music analysis \citep{Montesinos2020Solos}. More recent resources such as Saraga Audiovisual and PianoVAM further enrich the space with culturally specific concert material, hand landmarks, fingering information, and MIDI-aligned piano data \citep{Shankar2024SaragaAV,Kim2025PianoVAM}. On the modeling side, The Sound of Pixels and Music Gesture for Visual Sound Separation show that visible performer motion is informative for sound production and can guide source separation or audio-visual disentanglement \citep{Zhao2018SoundOfPixels,Gan2020MusicGesture}. Related multimodal classification work in other domains has also explored compact fusion pipelines, for example by combining lightweight detection-oriented backbones with multimodal decision modules \citep{Poloju2024MultimodalFusion}. This category is closest to our task because it preserves the public-video setting and treats body or hand motion as a first-class cue. Yet the literature remains fragmented across dataset construction, source separation, transcription, and specialized instrument settings, and most papers do not report a common protocol that jointly evaluates classification quality, computational cost, and failure under offset or corruption.

Taken together, existing methods have advanced multimodal music research from representation alignment to reasoning-rich understanding and then to performance-centered data and motion-aware analysis, but they still do not provide a unified setting for lightweight real-time instrument classification on public videos under local audio-visual mismatch. \method{} is positioned in this gap: it retains the multimodal perspective of prior work and the performance-centric emphasis of public video datasets, while framing the problem as instrument recognition with explicit phase-aware diagnostics rather than as a fully supervised performance-element benchmark.
